\documentclass[11pt,a4paper]{article}
\usepackage{fontspec}
\usepackage[danish]{babel}
\usepackage[a4paper,margin=23mm,headheight=15pt,footskip=18mm]{geometry}
\usepackage{amsmath,amssymb,mathtools}
\usepackage{enumitem}
\usepackage{xcolor}
\usepackage{xurl}
\usepackage{hyperref}
\usepackage{fancyhdr}
\usepackage{titlesec}
\usepackage{microtype}
\usepackage{array}
\usepackage{bookmark}

\setmainfont{DejaVu Serif}
\setsansfont{DejaVu Sans}
\setmonofont{DejaVu Sans Mono}

\definecolor{TitleBlue}{HTML}{17233A}
\definecolor{AccentBlue}{HTML}{355F96}
\definecolor{PaleBlue}{HTML}{EEF4FA}
\hypersetup{
  unicode=true,
  colorlinks=true,
  linkcolor=AccentBlue,
  urlcolor=AccentBlue,
  pdftitle={Halvperiodesymmetrier for følger modulo 3},
  pdfsubject={Matematisk formalisering og algoritmisk audit},
  pdfkeywords={recurrence, modulo 3, antiperiodicity, Fourier, companion matrix},
  pdfcreator={XeLaTeX}
}

\titleformat{\section}{\large\bfseries\color{TitleBlue}}{\thesection.}{0.55em}{}
\titleformat{\subsection}{\normalsize\bfseries\color{TitleBlue}}{\thesubsection}{0.55em}{}
\titlespacing*{\section}{0pt}{1.2ex plus .2ex}{.55ex}
\titlespacing*{\subsection}{0pt}{.9ex plus .2ex}{.35ex}
\setlist[itemize]{leftmargin=1.7em,itemsep=.20em,topsep=.35em}
\setlist[enumerate]{leftmargin=1.9em,itemsep=.20em,topsep=.35em}
\setlength{\parindent}{0pt}
\setlength{\parskip}{.48em}
\emergencystretch=2.5em
\setcounter{tocdepth}{2}

\pagestyle{fancy}
\fancyhf{}
\fancyhead[L]{\small\color{AccentBlue}Halvperiodesymmetrier modulo 3}
\fancyhead[R]{\small\color{AccentBlue}Formalisering og audit}
\fancyfoot[C]{\small\color{TitleBlue}\thepage/6}
\renewcommand{\headrulewidth}{0.35pt}
\renewcommand{\footrulewidth}{0pt}

\newcommand{\F}{\mathbb{F}}
\newcommand{\Z}{\mathbb{Z}}
\newcommand{\oplusword}{\mathbin{\Vert}}
\newcommand{\centerlift}[1]{\widetilde{#1}}

\begin{document}
\sloppy
\thispagestyle{empty}
\vspace*{20mm}
\noindent\noindent{\color{AccentBlue}\rule{\textwidth}{0.8pt}}\\[22mm]
{\raggedright\fontsize{24}{29}\selectfont\bfseries\color{TitleBlue}Halvperiodesymmetrier for følger\\ modulo 3\par}
\vspace{9mm}
{\fontsize{15}{20}\selectfont\color{AccentBlue}Matematisk formalisering og algoritmisk audit\par}
\vspace{17mm}
\noindent\fcolorbox{AccentBlue}{PaleBlue}{%
  \begin{minipage}{0.92\textwidth}
  \vspace{2mm}\small Dette notat adskiller stringent tre fænomener, som ofte blandes sammen: negation af cykler, halvperiode-antiperiodicitet og omvending. Det beviser det entydige modsatte skift for et ikke-nul primitivt ord, giver kombinatoriske, spektrale og matrixbaserede karakteriseringer og forbinder resultaterne med en udtømmende analysator for modulære rekurrenser. \vspace{2mm}
  \end{minipage}}
\vfill
{\color{TitleBlue}\small
Kodelager: \textsf{Modular-recurrence-symetry-analyzer-}\\
Samlet PDF-udgave - 5. august 2026
}\par
\vspace{18mm}
\noindent{\color{AccentBlue}\rule{\textwidth}{0.8pt}}
\clearpage
\setcounter{page}{1}

\tableofcontents
\clearpage


\section{Formål}
Dette notat formaliserer de fænomener, som i kodelageret \nolinkurl{59200/k-bonacci-pisano-finder} beskrives med de uformelle udtryk ``perioder komplementære til tre'' og ``selvkomplementær, når den deles i to''.

Den invariante ramme er periodiske ord over legemet
\[
\F_3=\{0,1,2\},
\]
med den additive involution
\[
\nu(x)=-x\pmod 3,\qquad \nu(0)=0,\quad \nu(1)=2,\quad \nu(2)=1.
\]
Udtrykket ``komplement til 3'' bør derfor erstattes af \emph{additiv invers modulo 3} eller \emph{komplement ved negation modulo 3}. Den decimale sum af de valgte repræsentanter er ikke altid 3, fordi 0 sendes på 0.

\section{Tre forskellige egenskaber}
Lad $w=(w_0,\ldots,w_{T-1})\in\F_3^T$ være et ord med primitiv periode $T$.

\subsection{Par af modsatte perioder}
To perioder $w$ og $v$ danner et modsat par, hvis
\[
v_n=-w_n\pmod 3
\]
for alle $n$. De kan tilhøre to forskellige cykler.

\subsection{Halvperiode-antiperiodicitet}
Ordet $w$ er halvperiode-antiperiodisk, hvis $T=2h$ og
\[
w_{n+h}=-w_n\pmod 3
\]
for alle $n$. Det kan da skrives
\[
w=H\oplusword(-H),\qquad H=(w_0,\ldots,w_{h-1}).
\]
Eksempler:
\[
01120221=0112\oplusword0221,\quad 011022=011\oplusword022,\quad 01220211=0122\oplusword0211.
\]

\subsection{Omvending af anden halvdel}
Operationen
\[
R(a_0,\ldots,a_{h-1})=(a_{h-1},\ldots,a_0)
\]
er uafhængig af negation. I Fibonacci-eksemplet er
\[
R(0221)=1220.
\]
Således er 1220 det omvendte antipodale ord til 0112, ikke selve den anden halvdel.

\clearpage
\section{Sætningen om det entydige modsatte skift}
\textbf{Sætning.} Lad $w$ være et ikke-nul primitivt ord med periode $T$. Hvis der findes et skift $r$, så
\[
w_{n+r}=-w_n
\]
for alle $n$, da er $T$ lige og
\[
r\equiv \frac{T}{2}\pmod T.
\]

\textbf{Bevis.} Anvendes skiftet to gange, fås $w_{n+2r}=w_n$. Primitivitet medfører $T\mid 2r$. Skiftet $r$ er ikke nul modulo $T$, fordi et ikke-nul ord over $\F_3$ ikke kan opfylde $w=-w$. Det entydige ikke-trivielle element af orden 2 i den cykliske gruppe $\Z/T\Z$ findes kun, når $T$ er lige, og er da $T/2$. \hfill$\square$

\textbf{Konsekvens.} En ikke-nul primitiv periode er derfor enten:
\begin{enumerate}
\item sendt ved negation til en anden cykel;
\item invariant under negation og dermed nødvendigvis halvperiode-antiperiodisk.
\end{enumerate}
Denne dikotomi fjerner forvekslingen mellem parret
\[
00111201,\qquad 00222102
\]
og intrinsisk antiperiodiske perioder som 011022.

\section{Kombinatoriske konsekvenser}
Hvis $w=H\oplusword(-H)$, så
\[
\#\{n:w_n=1\}=\#\{n:w_n=2\},\qquad \#\{n:w_n=0\}\equiv0\pmod2,
\]
og
\[
\sum_{n=0}^{T-1}w_n=0\quad\text{i }\F_3.
\]
Disse betingelser er nødvendige, men ikke tilstrækkelige. Med det centrerede løft
\[
\centerlift0=0,\qquad \centerlift1=1,\qquad \centerlift2=-1
\]
faktoriserer det genererende polynomium som
\[
W(X)=\sum_{n=0}^{2h-1}\centerlift{w_n}X^n=(1-X^h)\sum_{n=0}^{h-1}\centerlift{w_n}X^n.
\]

\clearpage
\section{Fourier-karakterisering}
Lad
\[
\widehat w(k)=\sum_{n=0}^{2h-1}\centerlift{w_n}e^{-2\pi i kn/(2h)}.
\]
Da gælder $w_{n+h}=-w_n$ præcis, når $\widehat w(k)=0$ for alle lige indekser $k$. Når $w_{n+h}=-w_n$,
\[
\widehat w(k)=\bigl(1-(-1)^k\bigr)\sum_{n=0}^{h-1}\centerlift{w_n}e^{-2\pi i kn/(2h)},
\]
så faktoren forsvinder for lige $k$. Omvendt ligger ordet i underrummet spændt af de ulige moder, hvis alle lige moder forsvinder. Et skift med $h$ virker på mode $k$ som $(-1)^k=-1$ og derfor som $-I$ på ordet. Dette er en eksakt spektral diagnostik, ikke blot en visuel sammenligning af halvdelene.

\section{Lineære følger modulo 3}
Betragt en rekurrens af orden $k$,
\[
u_{n+k}=c_0u_n+c_1u_{n+1}+\cdots+c_{k-1}u_{n+k-1}\pmod3.
\]
Tilstanden er $s_n=(u_n,\ldots,u_{n+k-1})^T$, og udviklingen er $s_{n+1}=Ms_n$, hvor $M$ er ledsagematricen.

\subsection{Ren periodicitet}
Determinanten af $M$ er, bortset fra fortegn, $c_0$. Over $\F_3$ gælder
\[
M\text{ er invertibel}\iff c_0\neq0.
\]
I så fald er enhver tilstandsbane rent periodisk. Hvis $c_0=0$, kan der forekomme en forperiode, som skal adskilles fra cyklen.

\subsection{Negation af cykler}
Linearitet giver $M(-s)=-M(s)$. Negation sender derfor hver cykel til en cykel af samme længde. For en ikke-nul cykel er kun to tilfælde mulige:
\begin{itemize}
\item to forskellige cykler udveksles ved negation;
\item én invariant cykel, nødvendigvis halvperiode-antiperiodisk.
\end{itemize}
Nulcyklen er den eneste degenererede undtagelse: den er punktvist fikseret af negation og definerer ikke en ikke-triviel primitiv antiperiode.

\subsection{Kriterium for én begyndelsestilstand}
For en begyndelsestilstand $s_0$ medfører relationen
\[
M^hs_0=-s_0
\]
for enhver lineær udlæsning af tilstanden, at
\[
u_{n+h}=-u_n.
\]
Når $s_0\neq0$, og cyklen er primitiv, giver relationen den ovenfor beskrevne halvperiode-antiperiode.

\clearpage
\subsection{Uniformt kriterium}
Alle begyndelsestilstande opfylder den samme antiperioderelation $h$, hvis og kun hvis
\[
M^h=-I.
\]
Ækvivalent opfylder minimalpolynomiet $\mu_M$
\[
\mu_M(X)\mid X^h+1\qquad\text{i }\F_3[X],
\]
og derfor er $M^{2h}=I$.

\section{Fibonacci-eksempel}
For
\[
M=\begin{pmatrix}0&1\\1&1\end{pmatrix}\quad\text{over }\F_3
\]
fås $M^4=-I$ og $M^8=I$. Perioden fra begyndelsestilstanden $(0,1)$ er
\[
01120221=0112\oplusword(-0112).
\]
Den omvendte anden halvdel er $R(0221)=1220$. Den numeriske sammenfald, ved decimal læsning,
\[
1220=L_{14}+F_{14}=2F_{15}
\]
er en ekstra kodningsidentitet, ikke en generel følge af antiperiodicitet.

\section{Generalisering til vilkårlige følger modulo 3}
Ingen rekurrensantagelse er nødvendig for at analysere et givet periodisk ord. For hver primitiv periode kan man beregne:
\begin{itemize}
\item dens rotationscykel;
\item dens additive inverse;
\item dens affine symmetrier $w_{n+r}=aw_n+b$;
\item dens omvendingssymmetrier $w_{r-n}=aw_n+b$;
\item enhver antiperiodicitet;
\item dens lige og ulige Fourier-moder.
\end{itemize}
Matrixteori er kun nødvendig, når en erklæret lineær rekurrens foreligger.

\section{Udvidelse til et generelt modulus}
Over $\Z/m\Z$ er den naturlige involution stadig $x\mapsto-x$. Definitionen
\[
w_{n+h}=-w_n\pmod m
\]
forbliver gyldig. For $m=2$ er negation identiteten, og antiperiodicitet falder sammen med almindelig periodicitet. For $m>2$ er forskellen ikke-triviel. Udtrykket ``komplement til $m$'' bør ikke bruges som algebraisk definition, fordi det afhænger af de valgte heltalsrepræsentanter; den invariante term er ``additiv invers modulo $m$''.

\clearpage
\section{Algoritmisk audit}
En periode for en rekurrens af orden $k$ modulo $m$ skal detekteres i det endelige tilstandsrum
\[
(\Z/m\Z)^k,
\]
ikke ved at lede efter en gentaget blok i et vilkårligt langt præfiks. Det medfølgende script bruger:
\begin{itemize}
\item kompakt base-$m$-kodning af tilstande;
\item Brents algoritme for én begyndelsestilstand, med $O(\mu+\lambda)$ tid og $O(1)$ hukommelse;
\item eksakt dekomposition af funktionsgrafen for alle tilstande, med $O(m^k)$ tid;
\item parallelisering på tværs af koefficientfamilier;
\item separat klassifikation af modsatte cykler og antiperiodiske cykler;
\item en matrixkontrol af $M^h=-I$.
\end{itemize}
Der kræves ingen Mathematica-kerne eller ekstern wrapper.

\section{Reproducerbare scanningsresultater}
Scanningen af de femten navngivne familier i kodelageret genfinder for tritetranacci-familien forskellige modsatte par med længder 24 og 8 samt de antiperiodiske perioder 011022, 01220211 og 12.

En anden scanning dækker de 119 ikke-nul binære koefficientvektorer af orden 2 til 6:
\[
\sum_{k=2}^{6}(2^k-1)=119.
\]
Scriptet finder:
\begin{itemize}
\item 13 familier, der opfylder det globale kriterium $M^h=-I$;
\item 88 familier med mindst én ikke-triviel halv-antiperiodisk cykel;
\item 96 familier med mindst ét særskilt par af modsatte cykler.
\end{itemize}
Dette er udtømmende eksperimentelle resultater inden for dette endelige vindue. De kan reproduceres fra \nolinkurl{mod3_symmetry_scan_summary.json}, \nolinkurl{binary_recurrence_families_mod3.json} og \nolinkurl{binary_recurrence_families_mod3_summary.csv}.

\section*{Konklusion}
Negation modulo 3 virker som en strukturel involution på periodiske ord og på cykler af lineære rekurrenser. For et ikke-nul primitivt ord kan invarians under denne involution kun forekomme gennem det entydige halvperiodeskift. Den samme egenskab genkendes ækvivalent ved faktorisering af det genererende polynomium, ved at de lige Fourier-moder forsvinder, eller, i den lineære ramme, ved matrixrelationen $M^h=-I$. Den algoritmiske analyse adskiller derfor intrinsiske cykelsymmetrier fra blotte forbindelser mellem modsatte cykler.


\end{document}
